\section{Introduction}

Machine learning has fundamentally transformed software development, but deploying ML systems to production 
remains a significant engineering challenge. The gap between successful model development and reliable 
production deployment is often underestimated in academic settings. This report documents a complete, 
end-to-end machine learning product: a movie rating prediction system that moves beyond model training 
to address the full lifecycle of an ML system.

\subsection{Motivation}

The movie recommendation domain serves as an excellent case study for several reasons:

\begin{enumerate}
  \item \textbf{Clear business value}: Movie rating predictions directly support recommendation engines, 
  which are critical for user engagement and platform success.
  
  \item \textbf{Rich collaborative patterns}: User-movie interactions exhibit strong collaborative signal, 
  where similar users rate similar movies similarly. This structure makes collaborative filtering effective.
  
  \item \textbf{Tractable scale}: The MovieLens 100K dataset is large enough to demonstrate real-world 
  challenges (sparsity, computation) while remaining manageable for classroom implementation.
  
  \item \textbf{Accessible baselines}: Collaborative filtering is well-understood and has established 
  theoretical foundations, allowing focus on engineering best practices rather than algorithmic novelty.
\end{enumerate}

\subsection{Scope of Work}

This project addresses three major engineering domains:

\textbf{Machine Learning:} We train a Singular Value Decomposition (SVD) model on the MovieLens 100K dataset 
to perform collaborative filtering. The model learns latent factor representations of users and movies, 
enabling rating predictions for arbitrary user-movie pairs.

\textbf{API and Systems Design:} We implement a RESTful API using FastAPI, providing clean interfaces 
for health monitoring and prediction serving. The API enforces input validation, handles errors gracefully, 
and maintains separation of concerns between model logic and request handling.

\textbf{Operations and Deployment:} We containerize the system using Docker and Docker Compose, ensuring 
portable and reproducible deployment. Health checks, environment configuration, and volume management 
demonstrate production-ready practices.

\subsection{Key Contributions}

This work makes the following contributions:

\begin{enumerate}
  \item A complete, working implementation of an ML system suitable for production-like constraints.
  
  \item Practical demonstration of API design patterns (validation, error handling, versioning) for ML systems.
  
  \item Comprehensive testing strategy that balances functional validation with realistic edge case handling.
  
  \item Dockerized deployment with health checks, supporting both local and containerized execution.
  
  \item Documentation and code structure that prioritizes maintainability and team collaboration.
\end{enumerate}

\subsection{Report Structure}

The remainder of this report is organized as follows:

\begin{itemize}
  \item \textbf{Section~\ref{sec:problem}}: Formal problem definition and rating prediction task formulation.
  
  \item \textbf{Section~\ref{sec:dataset}}: MovieLens 100K dataset characteristics, statistics, and preprocessing.
  
  \item \textbf{Section~\ref{sec:methodology}}: Theoretical foundations of collaborative filtering and matrix factorization.
  
  \item \textbf{Section~\ref{sec:model}}: Mathematical treatment of SVD and the Surprise library implementation.
  
  \item \textbf{Section~\ref{sec:architecture}}: System architecture covering training, artifact management, and serving.
  
  \item \textbf{Section~\ref{sec:api}}: RESTful API design, schema definitions, and error handling.
  
  \item \textbf{Section~\ref{sec:experiments}}: Training results, model evaluation, and prediction behavior analysis.
  
  \item \textbf{Section~\ref{sec:deployment}}: Docker containerization and deployment configuration.
  
  \item \textbf{Section~\ref{sec:testing}}: Testing strategy, test cases, and validation approach.
  
  \item \textbf{Section~\ref{sec:limitations}}: System limitations, failure modes, and future improvements.
  
  \item \textbf{Section~\ref{sec:conclusion}}: Summary and directions for future work.
\end{itemize}

\vspace{0.5cm}

This report emphasizes the engineering and operational aspects of ML systems, recognizing that 
successful ML in production requires far more than accurate model training.
